\section{Materials and Methods}
\label{sec:methods}
This section defines the exact image decomposition, the twin lesion-invariance objective and the counterfactual audit
(Section~\ref{sec:m-formulation}) and describes the data, their preprocessing and splits (Section~\ref{sec:m-data}). It then describes
the suppressors of Study~I and their evaluation (Section~\ref{sec:m-study1}), the downstream models (Section~\ref{sec:m-downstream}), the
fusion study of Study~II (Section~\ref{sec:m-study2}), the external validation (Section~\ref{sec:m-bdcxr}) and the statistical analysis
(Section~\ref{sec:m-stats}). Algorithm~\ref{alg:pipeline} summarises the pipeline and Fig.~\ref{fig:architecture} shows the models; the
rendering of the first-generation (v1) training images from LIDC-IDRI CTs is described in Appendix~\ref{app:v1}. Table~\ref{tab:notation}
lists the notation.

\begin{table}[!t]
\centering
\caption{Notation.}
\label{tab:notation}
\footnotesize
\renewcommand{\arraystretch}{1.08}
\setlength{\tabcolsep}{3pt}
\begin{tabular}{@{}>{\raggedright\arraybackslash}p{0.25\linewidth}>{\raggedright\arraybackslash}p{0.71\linewidth}@{}}
\toprule
Symbol & Meaning \\
\midrule
$X$, $S$, $B$ & radiograph, soft-tissue and bone image of a training sample, $S=X-B$ \\
$x$, $b$ & suppressor input and its exact bone target \\
$f$, $f_k$ & suppressor (predicts the bone image); ensemble member $k$ \\
$\bar f$, $\sigma$ & ensemble mean and per-pixel standard deviation ($K=3$) \\
$\hat s$, $\hat B$ & soft-tissue estimate $x-\bar f(x)$; bone image $X-S$ given to the classifier \\
$n$, $\alpha$ & inserted lesion signal and its amplitude \\
$x^{+}$, $M$ & twin $\min(x+n,1)$; lesion support dilated by 3\,px \\
$w$ & twin weight ($w=16$ in the final model) \\
$A$ & lesion absorption, Eq.~\eqref{eq:absorption} \\
$q$, $\mathcal{C}_p$ & conformal scale and interval at pixel $p$ \\
$s$, PSNR$_{\mathrm{cal}}$ & dual-energy soft-tissue reference; calibrated PSNR \\
$R_{\mathrm{rib}}$, $R_{\mathrm{non}}$ & gradient-energy reduction on rib-boundary and other lung pixels \\
$\Delta_G$, $\rho_{c,s}$ & mean AUROC difference of a finding group $G$; routing of class $c$ to stream $s$ \\
$I$ & pointing-game interaction, focal minus other targets \\
\bottomrule
\end{tabular}
\end{table}

\begin{figure*}[!t]
\centering
\safefig[0.9\textwidth]{figures/fig_tmi_architecture.pdf}
\caption{Model architectures. (a) Study~I: bone suppressor trained with the twin lesion-invariance regulariser
(Section~\ref{sec:m-twin}); the twin $x^{+}=\min(x+n,1)$ shares the weights and the bone target of $x$. (b) Study~II: multi-stream
classifier with a shared encoder and a class-query decoder (Section~\ref{sec:m-study2}), and the five input configurations.}
\label{fig:architecture}
\end{figure*}

\begin{algorithm}[!t]
\caption{Training, counterfactual audit and downstream evaluation}
\label{alg:pipeline}
\footnotesize
\begin{algorithmic}[1]
\REQUIRE Qure.ai projections $\mathcal{D}_{\mathrm{qp}}$ with exact $(X,B)$ (Eq.~\eqref{eq:qure}); lesion libraries
 $\mathcal{L}_{\mathrm{train}}$, $\mathcal{L}_{\mathrm{test}}$ from LIDC-IDRI training and test CTs; real test sets
\ENSURE Suppressor ensemble $\bar f$ with intervals $\mathcal{C}_p$; downstream and fusion classifiers; all estimates with 95\,\% intervals
\FOR[\textit{Study~I}]{seed $k\in\{0,1,2\}$}
 \FOR{step $=1$ \TO \num{37500}}
  \STATE Compose a batch: re-weight $B$, shared affine, insertion, tone curve $\gamma$, $B'=X'-S'$ (Eq.~\eqref{eq:tone})
  \STATE Build twins $x^{+}=\min(x+n,1)$, $n$ from $\mathcal{L}_{\mathrm{train}}$, for half of the batch
  \STATE Update $f_k$ with AdamW on $\mathcal{L}$ of Eq.~\eqref{eq:twin}; update EMA weights
 \ENDFOR
\ENDFOR
\STATE $\bar f,\sigma\gets$ Eq.~\eqref{eq:ensemble}; calibrate $q$ on Qure.ai validation pixels (Eq.~\eqref{eq:conformal})
\FOR[\textit{audit}]{each real test image $x$}
 \STATE Draw six lung sites with a fixed seed and lesions from $\mathcal{L}_{\mathrm{test}}$
 \STATE For every model, compute $A$ (Eq.~\eqref{eq:absorption}) and bone overlap from independent masks
\ENDFOR
\STATE Score fidelity, rib benchmark, $\Delta$CNR, detection and classification; test H1$'$--H5$'$ and E1--E3
\FOR[\textit{Study~II}]{config $\in\{$orig, soft, early, hybrid, ctrl$\}$, data set, fold, seed $0..4$}
 \STATE Train the class-query classifier on the $(X,S,\hat B)$ streams (Eqs.~\eqref{eq:tokens}--\eqref{eq:bce})
\ENDFOR
\STATE Seed ensembles (Eq.~\eqref{eq:seedens}); cluster-bootstrap contrasts with Holm; test H6--H10; pointing game, conformal sets, AURC
\RETURN all estimates with 95\,\% intervals, whatever their direction
\end{algorithmic}
\end{algorithm}


% ---------------------------------------------------------------------------------------------------------------------
\subsection{Problem Formulation and Counterfactual Audit Definitions}
\label{sec:m-formulation}

\subsubsection{Training images with an exactly known bone image}
\label{sec:m-qure}
A suppressor $f$ predicts the bone image of a radiograph; the suppressed (soft-tissue) image is the radiograph minus that prediction.
Supervision needs training images whose bone image is known exactly. The final (v2) suppressors are trained on public projections of chest
CT \cite{angaitkar2026anatomy}: each row of the Qure.ai data holds a composite radiograph $T$ and the translated bone, soft-tissue and lung
projections $B_t$, $S_t$ and $L$ (8-bit, $512\times512$, one sample per CT volume at 100~kVp). We verified that
\begin{equation}
\begin{gathered}
T=a\,S_t+\beta\,B_t+c\,L+d,\\
X=\frac{T}{255},\quad B=\frac{\beta\,B_t}{255},\quad S=X-B
\end{gathered}
\label{eq:qure}
\end{equation}
holds up to rounding, with a per-sample bone weight $\beta$ fitted by least squares, so $B$ is the exactly additive bone image of the
normalised input $X$. Over all \num{21887} rows the least-squares fit leaves a mean absolute residual of 0.36 grey levels per projection
(largest 3.2 of 255).

A training sample is composed on the GPU in this order. The bone image is re-weighted ($\tilde B=w_bB$, $w_b$ log-uniform in $[0.6,1.7]$,
so $\tilde X=S+\tilde B$ stays exact); one random affine transform is applied to all maps; in the arms with counterfactual insertion
(Table~\ref{tab:arms}), up to three lesions are inserted into $S$ with the bone target untouched; and a monotone gamma tone curve $\gamma$
gives
\begin{equation}
X'=\gamma(\tilde X),\qquad S'=\gamma(S),\qquad B'=X'-S',
\label{eq:tone}
\end{equation}
which keeps the decomposition exact under any monotone mapping. Linear blur is applied to $X'$ and $B'$, and noise and scatter to $X'$
only. The first-generation (v1) suppressors were trained on radiographs rendered from LIDC-IDRI CTs with polychromatic attenuation physics
(Appendix~\ref{app:v1}).

\subsubsection{Suppression loss}
For a predicted bone image $\hat b=f(x)$ and the target $b$, the suppression loss is
\begin{equation}
\begin{split}
\mathcal{L}_{\mathrm{sup}}(\hat b,b)={}&0.5\,\lVert\hat b-b\rVert_1+0.5\,\lVert\hat b-b\rVert_2^2\\
&+0.5\,\big(1-\mathrm{SSIM}(\hat b,b)\big)\\
&+2.5\,\big(\lVert\hat b-b\rVert_1+\lVert\hat b-b\rVert_2^2\big)_{\mathrm{ROI}},
\end{split}
\label{eq:loss}
\end{equation}
with norms taken as means over pixels and the region of interest (ROI) taken from the lung component. The suppressed image is
$x-\hat b$.

\subsubsection{Twin lesion-invariance objective}
\label{sec:m-twin}
Let $f$ be the suppressor, $x$ a training image and $b$ its exact bone image. A \emph{twin} of $x$ is $x^{+}=\min(x+n,1)$, where $n\ge0$ is
the signal of a few lesion patches taken from the lesion library of \emph{training} CTs and pasted at random positions in the lung fields,
in composite units and after all augmentation. The bone target of $x^{+}$ is the same $b$. Let $M$ be the lesion support dilated by
$3$~px. The loss is
\begin{equation}
\begin{gathered}
\mathcal{L}=\tfrac12\big[\mathcal{L}_{\mathrm{sup}}(f(x),b)+\mathcal{L}_{\mathrm{sup}}(f(x^{+}),b)\big]+w\,\mathcal{L}_{\mathrm{inv}},\\
\mathcal{L}_{\mathrm{inv}}=\frac{\sum_p M_p\,\big|f(x^{+})_p-f(x)_p\big|}{\max\big(1,\sum_p M_p\big)},
\end{gathered}
\label{eq:twin}
\end{equation}
where $w$ is the twin weight. Half of every batch consists of twins, so the cost per step does not change. The supervised term already asks
for the unchanged bone image on $x^{+}$; $\mathcal{L}_{\mathrm{inv}}$ ties the two predictions directly on the lesion support. Each twin
carries one or two pasted lesions (uniform), each with a random scale and amplitude; the supervised term is averaged over all base images
and twins of the batch, and $\mathcal{L}_{\mathrm{inv}}$ is normalised per image by its support and then averaged over the twins.
Fig.~\ref{fig:insertion} shows one pasted lesion, its support $M$ and the unchanged bone target.

\begin{figure}[!t]
\centering
\safefig[\linewidth]{figures/fig_lesion_insertion.pdf}
\caption{Counterfactual lesion insertion, illustrated on a rendered LIDC-IDRI test CT; in training, the twins of the regulariser are built
in the same way from Qure.ai projections with lesions of the training library. (a) Lesion signal $n$ of a 15~mm LIDC-IDRI nodule; (b) the
composite with the lesion, $x+n$, and the dilated support $M$; (c) the bone target under the site, which the insertion leaves unchanged.}
\label{fig:insertion}
\end{figure}

\subsubsection{Counterfactual lesion absorption}
\label{sec:m-absorption}
For a suppressor $f$ that outputs the bone image, an image $x$ and a lesion signal $n\ge0$ in composite-image units, the absorption at
amplitude $\alpha$ is
\begin{equation}
A=\frac{\big\langle f(x+\alpha n)-f(x),\;\alpha n\big\rangle}{\big\langle \alpha n,\;\alpha n\big\rangle},
\label{eq:absorption}
\end{equation}
with both inner products taken over the lesion support (the pixels where the inserted signal exceeds 2\,\% of its peak). $A=0$ means that
the lesion stays in the soft-tissue image and $A=1$ that it was removed with the bone; $A$ is a regression slope and can fall outside
$[0,1]$. A model that does nothing has $A=0$, so $A$ is reported together with a suppression gain and never alone. Two suppression
measures are paired with $A$: the bone-region PSNR on the held-out SZCH pairs and, on VinDr-RibCXR, the reduction of gradient energy on
rib-boundary pixels (the rib benchmark, Eq.~\eqref{eq:rib}).

% ---------------------------------------------------------------------------------------------------------------------
\subsection{Data Sets, Preprocessing and Splits}
\label{sec:m-data}
Ten data sets serve three roles in this work: sources for training the suppressors, real radiographs for evaluating them and the
downstream models, and an external paediatric set for testing transfer to an unseen population. Table~\ref{tab:data} lists each set with
its role, splits and licence or access condition, Fig.~\ref{fig:data_splits} shows the size of every split, and
Fig.~\ref{fig:data_pipeline} with Table~\ref{tab:preproc} follows each image source from its raw form to the study that uses it
(Section~\ref{sec:m-preproc}).

\subsubsection{Training sources}
The LIDC-IDRI collection \cite{armato2011lidc,armato2015lidctcia,clark2013tcia}, with some series read from the LUNA16 copy
\cite{setio2017luna16}, supplies the v1 renders, the lesion library used to paste lesions into training images, and the synthetic test CTs.
We used 1\,018 chest CT series from 1\,010 patients with the radiologists' nodule contours. Series were read from DICOM or from the LUNA16
MetaImage copy (voxel-identical on all 12 series we compared, including nodule slice locations); 1\,000 series could be rendered (18 were
skipped because of non-axial orientation or irregular slice spacing). A render was \emph{usable} if the lungs were not cut by the scan, at
most 50 of the 512 frame rows lay beyond the scan, the lung volume was between 1.5 and 9~L and the bone volume exceeded 300~cm$^3$ (844
usable CTs). Patients were assigned to train/validation/test (70/10/20\,\%) by a hash of the patient identifier (588/98/158 CTs;
1\,057/172/298 consensus nodules) before any model was trained.

The v2 models are trained on the public Qure.ai \texttt{ct2xr-projections} data \cite{angaitkar2026anatomy}, which hold, for each CT
volume, a composite image and translated bone, soft-tissue and lung projections (Section~\ref{sec:m-qure}). The release holds \num{21887}
projections, one per CT volume; the volume names are those of CT-RATE \cite{hamamci2026ctrate}. We split them by patient (\num{14374}
patients) into \num{21051} training, 406 validation and 430 test projections.

\subsubsection{Real radiographs}
\emph{SZCH-X-Rays} holds paired radiographs and dual-energy soft-tissue images from one partner hospital \cite{sun2026bsldm,sun2026debonedit}.
We use the 75-pair test subset; its 75 images are split into a development part (24 images) and a held-out part (51 images) by a hash of
the file name. The 75 images are the published test split of the 741-pair version of the data set (592/74/75), acquired with dual-exposure
DES on a GE Discovery XR656 \cite{sun2026debonedit}; an earlier version of 818 pairs was used for BS-LDM \cite{sun2026bsldm}, and the
papers do not explain the difference.

\emph{VinDr-RibCXR} holds 245 radiographs with per-rib masks \cite{nguyen2021ribcxr}. The data are licensed and are processed only on the
licence holder's own computer and private cloud account; only metrics leave them.

\emph{Nodule and detection sets.} JSRT \cite{shiraishi2000jsrt} (154 nodules), Node21 \cite{sogancioglu2024node21} (1\,476 nodule boxes, of
which two narrower than 4\,px at 512\,px have no defined contrast, so the contrast analysis uses 1\,474), TBX11K \cite{liu2020tbx11k} and
VinBigData (the Kaggle competition data, related to VinDr-CXR \cite{nguyen2022vindrcxr}) are used for real-lesion contrast and for
downstream detection. The detection sets are split by a fixed hash of the image identifier: Node21 \num{3440}/\num{512}/\num{930}, TBX11K
\num{6200}/\num{688}/\num{2088} (its official validation set serves as test set) and VinBigData \num{10433}/\num{1469}/\num{2943} images
for training, validation and test. VinBigData boxes are kept where at least two radiologists marked the same finding; images with only
single-reader findings are left out. No image of any set was unreadable. VinBigData is the Kaggle competition built on the VinDr-CXR images
\cite{vinbigdata2020kaggle}; its rules allow non-commercial research use.

\emph{NIH ChestX-ray14} \cite{wang2017chestxray} is used for classification with its official test list (\num{25596} images) and a
patient-level split of the remaining images into \num{77890} training and \num{8634} validation images. The downstream models of Study~I
and the fusion study of Study~II use the same NIH and VinBigData splits; on VinBigData an image is positive for one of the 14 findings when
at least two of its readers marked it. The 984 expert boxes on 880 NIH test images (eight classes) are used only to score attribution maps.

\subsubsection{External validation set and near-duplicate screening}
\emph{BDCXR-3257} holds 3,257 paediatric chest radiographs (JPEG) from the Radiology and Imaging Department of Dr.~M~R~Khan Shishu Hospital
\& Institute of Child Health, Dhaka, collected in routine care and anonymised, with one label each: normal (880), bacterial pneumonia
(1,477) or viral pneumonia (900) \cite{durjoy2026bdcxr}, a task studied before on other paediatric sets \cite{kermany2018identifying}. Most
children are under five years old (1,132 images under one year). The data set is public on figshare under CC BY 4.0, but we obtained it
under the ethical approval REC-FSIT-2025/No:~11928 (Daffodil International University) and a data use agreement between the hospital and
the authors' research group, whose terms we follow: no radiograph of this set is shown, and no image or model derived from it is
redistributed by us. There are no patient identifiers, so near-duplicate images are grouped by their correlation at $64$~px
($r\ge0.93$; 48 pairs, the rest single images), and a group never spans two folds.

\subsubsection{Development and locked sets}
The locked test sets were fixed before any locked-set evaluation. After the ablation of the v1 suppressors the SZCH pairs had been looked
at, so SZCH became a development set for decisions about the final recipe, which was fixed before any held-out or locked evaluation of a
final model. Decisions used the 24 development images, the synthetic validation data and the real-image audit on all 75 SZCH images; the
audit on SZCH is therefore development evidence only, and the confirmatory audit runs on VinDr-RibCXR (not opened by any model run before
the final protocol was written) and JSRT. Paired fidelity claims use the 51 held-out SZCH pairs. While the recipe was being chosen,
screening tools printed held-out PSNR and SSIM next to the development numbers for v1-family and screening runs; no rule, table or decision
used them, and the tool was changed to print development numbers only.

\subsubsection{Preprocessing steps}
\label{sec:m-preproc}
Every image passes through a fixed chain before any model sees it (Fig.~\ref{fig:data_pipeline}, Table~\ref{tab:preproc}):
\begin{enumerate}
\item \emph{Decomposition of CT into exact components.} LIDC-IDRI CTs are split into bone, non-lung soft tissue and lung by Hounsfield
 rules, projected along the anteroposterior axis and rendered with polychromatic physics (Appendix~\ref{app:v1}); the Qure.ai projections
 are converted into an exactly additive radiograph and bone image (Eq.~\eqref{eq:qure}).
\item \emph{Augmentation that keeps the decomposition exact.} Bone re-weighting, one shared affine transform, lesion insertion and a
 monotone tone curve are applied so that $B'=X'-S'$ holds for every training sample (Eq.~\eqref{eq:tone}).
\item \emph{Real radiographs.} Each radiograph is resized to 1024\,px by area averaging and min--max scaled; the suppressors read it at
 512\,px, the classifiers at 256\,px and the detector at 512\,px.
\item \emph{Lung masks and audit sites.} Lung fields come from a published segmentation network, eroded by 49\,px, and six lesion sites
 per image are drawn once with a fixed seed, so that every model is audited at the same places.
\item \emph{Integrity screening and splits.} Unusable renders are removed by fixed rules, near-duplicate paediatric images are grouped, and
 every split is made by patient, image hash or duplicate group before any model sees the corresponding test data.
\end{enumerate}

\begin{figure*}[!t]
\centering
\safefig[0.88\textwidth]{figures/fig_data_pipeline_tmi.pdf}
\caption{Data preprocessing. Each row follows one image source from its raw form (left) through preprocessing and splitting (train /
validation / test; dev: development pairs, test: held-out pairs) to the studies that use it (right; I, II). Thumbnails are real images
where the licence allows: a LIDC-IDRI CT slice \cite{armato2015lidctcia} (CC~BY~3.0), a SZCH-X-Rays radiograph with its dual-energy
soft-tissue image \cite{sun2026bsldm} (CC~BY-NC~4.0) and an NIH ChestX-ray14 radiograph \cite{wang2017chestxray}. $\dagger$: no image
derived from CT-RATE is shown, so the Qure.ai row shows the composite radiograph $X$ and the exact bone image $B$ of a rendered LIDC-IDRI
CT. $\ddagger$: dummy radiographs, variants (contrast, noise, mirroring, framing) of the same rendered LIDC-IDRI CT, stand in for data sets
whose images are not shown (data use agreements or restrictive licences: VinDr-RibCXR, JSRT, VinBigData, BDCXR-3257, Node21; TBX11K is
not shown either); they are not images of these data sets.}
\label{fig:data_pipeline}
\end{figure*}

\begin{figure}[!t]
\centering
\safefig[\columnwidth]{figures/fig_data_splits_tmi.pdf}
\caption{Size of every data split on a logarithmic scale. SZCH-X-Rays has only development and held-out pairs, VinDr-RibCXR and JSRT are
test sets only, and BDCXR-3257 is used in five-fold cross-validation over near-duplicate groups.}
\label{fig:data_splits}
\end{figure}

\begin{table*}[!t]
\centering
\caption{Data sets, their role, splits (train / validation / test unless stated) and licence or access condition, as stated on the
providers' pages (read on 3 November 2026).}
\label{tab:data}
\footnotesize
\renewcommand{\arraystretch}{1.08}
\begin{tabularx}{\textwidth}{@{}>{\raggedright\arraybackslash}p{3.1cm}>{\raggedright\arraybackslash}p{4.4cm}%
>{\raggedright\arraybackslash}p{4.2cm}>{\raggedright\arraybackslash}X@{}}
\toprule
Data set & Role & Split & Licence or access \\
\midrule
LIDC-IDRI (TCIA) & v1 renders, lesion library, synthetic test CTs & 588 / 98 / 158 CTs (by patient) & CC BY 3.0; attribution with the data DOI \\
Qure.ai ct2xr-projections & v2 training source (one projection per CT-RATE volume) & \num{21051} / 406 / 430 projections (by patient) & CC BY-NC-SA 4.0 \\
SZCH-X-Rays & paired DES evaluation, site calibration & 24 development / 51 held-out pairs & CC BY-NC 4.0 \\
VinDr-RibCXR & locked audit, rib benchmark & 245 test images & licensed; local use, metrics only \\
JSRT & audit, nodule contrast & 247 test images (154 nodules) & free after registration; no redistribution or commercial use \\
Node21 & nodule contrast, detection & \num{3440} / 512 / 930 images & CC BY-NC-ND 4.0 \\
TBX11K & detection & \num{6200} / 688 / \num{2088} images & CC BY 4.0 (data) \\
VinBigData / VinDr-CXR & detection; fusion study & \num{10433} / \num{1469} / \num{2943} images & competition rules: non-commercial research use \\
NIH ChestX-ray14 & downstream classification; fusion study & \num{77890} / \num{8634} / \num{25596} images (test: official list) & unrestricted use with attribution \\
BDCXR-3257 & external validation; fusion study & five folds over near-duplicate groups (\num{3257} images) & CC BY 4.0; used under the providers' data use agreement \\
\bottomrule
\end{tabularx}
\end{table*}

\begin{table*}[!t]
\centering
\caption{Preprocessing of each image source. All splits are by patient, by a fixed hash of the image identifier or, for BDCXR-3257, by
near-duplicate group, and were fixed before any model saw the corresponding test data.}
\label{tab:preproc}
\footnotesize
\renewcommand{\arraystretch}{1.1}
\setlength{\tabcolsep}{3pt}
\begin{tabularx}{\textwidth}{@{}>{\raggedright\arraybackslash}p{3.0cm}>{\raggedright\arraybackslash}p{3.0cm}%
>{\raggedright\arraybackslash}X>{\raggedright\arraybackslash}p{3.4cm}@{}}
\toprule
Source & Used by & Preprocessing & Integrity checks \\
\midrule
LIDC-IDRI CT & v1 suppressors, lesion library & Bone, non-lung soft tissue and lung masks; $\rho=\mathrm{clip}(1+\mathrm{HU}/1000,0,3)$; AP
 projection of each component; body crop, $512\times512$; physics of Eq.~\eqref{eq:physics} at 80--130~kVp (training) or 100~kVp
 (evaluation); composer of Eq.~\eqref{eq:compose} & 18 of 1\,018 series not renderable; usability rules leave 844 CTs; DICOM and LUNA16
 copies voxel-identical on 12 series\\
Qure.ai projections & v2 suppressors & 8-bit, $512\times512$; $X=T/255$, exact $B$ (Eq.~\eqref{eq:qure}); bone weight log-uniform in
 $[0.6,1.7]$, shared affine, insertion, gamma tone curve, blur on $X'$ and $B'$, noise and scatter on $X'$ & Mean absolute residual of the
 additive fit 0.36 grey levels (largest 3.2 of 255)\\
Real radiographs (suppressor input) & Studies I and II & Suppressor run at 512\,px; lung masks from the TorchXRayVision PSPNet at
 1024\,px, eroded by 49\,px for audit sites & Masks exported once and shared by all models\\
Classifier inputs & Study~II, downstream classification, BDCXR-3257 & $256\times256$\,px; per-image standardisation and a learned affine
 map per stream; rotation $\le7^\circ$, scale 0.9--1.1, translation $\le3$\,\%, brightness and contrast, identical across channels &
 VinBigData labels need two readers\\
Detector inputs & Downstream detection & $512$\,px; boxes of findings marked by two or more readers & Node21: two boxes narrower than 4\,px
 excluded from contrast\\
BDCXR-3257 & External validation & JPEG; five class-stratified folds over near-duplicate groups (fold seed 0) & Correlation at 64\,px,
 $r\ge0.93$: 48 pairs grouped\\
\bottomrule
\end{tabularx}
\end{table*}

% ---------------------------------------------------------------------------------------------------------------------
\subsection{Study~I: Lesion-Faithful Bone Suppression and Deep Ensemble Uncertainty}
\label{sec:m-study1}
Study~I asks whether a CT-supervised suppressor can remove bone without removing lesions, and whether a twin regulariser can teach it to.
This subsection describes the suppressors, their training arms and comparators, the ensemble and its intervals, the audit, and the
measures used to score them on real radiographs.

\subsubsection{Suppressor architecture and training}
The suppressor (Fig.~\ref{fig:architecture}a) is a nested (UNet++) decoder \cite{zhou2018unetpp} over the four stages of a pretrained
encoder, with CBAM attention \cite{woo2018cbam}, a plain U-Net tail with two shallow convolutional skips to full resolution and a
$1\times1$ head that outputs the bone image directly. The v1 suppressors use a SwinV2-Tiny encoder \cite{liu2022swinv2} (ImageNet-1k,
window 8, $512^2$ input); relative to the reference this is a smaller backbone at half the resolution, forced by compute. They were trained
with AdamW (learning rate $2\times10^{-4}$, weight decay 0.01, 500-step warm-up, cosine decay), batch 8, mixed precision, gradient clipping
at 1, EMA weights and \num{200000} samples per arm. The dose-response and ablation arms of v2 use the same SwinV2-T encoder; our final model
uses a ConvNeXt-V2-T encoder \cite{woo2023convnextv2} (FCMAE pre-training, ImageNet-22k then ImageNet-1k fine-tuning
\cite{deng2009imagenet}; \texttt{timm} \cite{wightman2019timm} model \texttt{convnextv2\_tiny.fcmae\_ft\_in22k\_in1k}) with a learning rate
$0.25\times$ that of the decoder. Every v2 run is trained at $512$~px with AdamW \cite{loshchilov2019adamw} (learning rate
$2\times10^{-4}$, batch 8), the schedule and loss above, and \num{300000} samples (\num{37500} steps); evaluation uses the EMA weights. The
encoder screen uses \num{100000} samples. Every arm changes exactly one factor of its parent (Table~\ref{tab:arms}); all other settings,
including the seed set, are identical. Table~\ref{tab:settings} collects the training settings of every model in this work.

\begin{table*}[!t]
\centering
\caption{Training settings of all models. Learning rates are those of AdamW unless stated; EMA: exponential moving average of the
weights.}
\label{tab:settings}
\footnotesize
\renewcommand{\arraystretch}{1.08}
\setlength{\tabcolsep}{3pt}
\begin{tabularx}{\textwidth}{@{}>{\raggedright\arraybackslash}p{2.6cm}>{\raggedright\arraybackslash}X%
>{\raggedright\arraybackslash}p{4.4cm}>{\raggedright\arraybackslash}p{2.5cm}>{\raggedright\arraybackslash}p{2.7cm}@{}}
\toprule
Model & Architecture and input & Optimisation & Budget and seeds & Selection \\
\midrule
v1 suppressors & SwinV2-T encoder, UNet++/CBAM decoder; LIDC-IDRI renders at 512\,px & $2\times10^{-4}$, weight decay 0.01, 500-step
 warm-up, cosine; batch 8; mixed precision, clipping at 1; EMA & \num{200000} samples; 3 seeds per arm & validation CTs \\
Final suppressors (ours) & ConvNeXt-V2-T encoder (learning rate $0.25\times$), UNet++/CBAM decoder; Qure.ai projections at 512\,px; twin
 weight $w=16$ & $2\times10^{-4}$, batch 8, same schedule; EMA & \num{300000} samples (\num{37500} steps); seeds 0--2, seed 3 reserve &
 development data, rule fixed in advance \\
Dose-response, ablation & SwinV2-T encoder, otherwise as ours & as ours & \num{300000} samples; 2--4 seeds per arm & -- \\
Encoder screen & ConvNeXt V2, DINOv3, RAD-DINO (learning rate $0.25\times$) & as ours & \num{100000} samples & screening rule \\
Downstream classifier & DenseNet-121, 256\,px & $3\times10^{-4}$, weight decay $10^{-4}$, batch 64, cosine & 10 epochs; 3 seeds & best
 validation macro-AUROC \\
Downstream detector & Faster R-CNN, ResNet-50 FPN, 512\,px & SGD 0.01, momentum 0.9, weight decay $10^{-4}$, batch 8 & 12 epochs; 3
 seeds & best validation CPM \\
Fusion classifier & shared DenseNet-121, class-query decoder (2 layers, 8 heads, $D=256$), 256\,px & $3\times10^{-4}$, weight decay
 $10^{-4}$, batch 32, 200 warm-up steps, cosine; mixed precision & 20 epochs; seeds 0--4 (each BDCXR-3257 fold) & best validation
 macro-AUROC \\
BDCXR-3257 classifier (E1) & as the downstream classifier & as the downstream classifier & 10 epochs; 3 seeds per fold & 10\,\% of each
 training part, by group \\
\bottomrule
\end{tabularx}
\end{table*}

\subsubsection{Training arms}
Table~\ref{tab:arms} lists the arms. The dose-response varies the twin weight over $w\in\{0,4,8,16,32,64\}$ in the SwinV2-T family with 3,
2, 2, 4, 2 and 2 seeds ($w=1$ appears only in the screening). The ablation changes one part of the regulariser at a time: a stop-gradient
on $f(x)$; a support dilated further with a $31\times31$ window, or the whole image, instead of $M$; up to four instead of up to two pasted
lesions; and lesion insertion into the base image without the tie term. The insertion-only arm and the last four rows of
Table~\ref{tab:arms} form this ablation (two seeds each, SwinV2-T, $w=16$, development data).

\begin{table}[!t]
\centering
\caption{Training arms. ``Lesions'' is the maximum number of lesions inserted into the base image (before the physics for the LIDC
renders, before the tone curve for the Qure.ai projections); the twin arms paste 1--2 further lesions per twin (up to 4 in
\texttt{qp\_twin16\_n4}) into the finished composite. ``Twin'' is the twin weight $w$, varied over $\{4,8,16,32,64\}$ in the
\texttt{qp\_twin}$w$ arms ($w=1$ in the screening only). The parent of an arm is \texttt{recipe} for the LIDC arms, \texttt{qp} for the
first Qure.ai arms and \texttt{qp\_twin16} for the last four rows.}
\label{tab:arms}
\footnotesize
\setlength{\tabcolsep}{3pt}
\begin{tabular}{@{}lccc>{\raggedright\arraybackslash}p{0.36\linewidth}@{}}
\toprule
Arm & Source & Lesions & Twin & Change \\
\midrule
\texttt{recipe} & LIDC & 0 & none & v1 baseline \\
\texttt{nodule\_loss} & LIDC & 0 & none & lesion-weighted loss \\
\texttt{clia} & LIDC & 3 & none & counterfactual insertion \\
\texttt{clia\_nodule} & LIDC & 3 & none & insertion and lesion-weighted loss \\
\texttt{qp} & Qure & 0 & none & Qure.ai projections (no-twin control) \\
\texttt{qp\_clia} & Qure & 3 & none & insertion only, no tie term \\
\texttt{qp\_twin}$w$ & Qure & 3 & varied & twin regulariser \\
\texttt{qp\_twin16\_sg} & Qure & 3 & 16 & stop-gradient on the un-pasted branch \\
\texttt{qp\_twin16\_wide} & Qure & 3 & 16 & support dilated with a $31\times31$ window \\
\texttt{qp\_twin16\_all} & Qure & 3 & 16 & whole-image support \\
\texttt{qp\_twin16\_n4} & Qure & 3 & 16 & four lesions per twin \\
\bottomrule
\end{tabular}
\end{table}

\subsubsection{Encoder screen and model selection}
\label{sec:m-encoders}
To check whether a newer encoder helps under the same recipe, ConvNeXt V2 \cite{woo2023convnextv2}, DINOv3 \cite{simeoni2025dinov3} and
RAD-DINO \cite{perezgarcia2025raddino} were trained at a reduced budget, and an encoder moved to the full budget only if it met a rule
fixed in advance: at \num{100000} samples, absorption $A\le0.070$ on the development audit with bone-region PSNR$_{\mathrm{cal}}\ge21.60$~dB
on the development pairs, or PSNR$_{\mathrm{cal}}\ge21.95$~dB with $A\le0.100$ (Supplementary Table~S7). The thresholds come from the
SwinV2-T range at that budget and were written after one SwinV2-B run had been seen. Pretrained encoders other than SwinV2 are trained with
a learning rate $0.25\times$ that of the decoder; ConvNeXt-V2-T diverged at the full rate (fp16 overflow) and was rerun at $0.25\times$, an
extension of the rule written for vision transformers. DINOv3 ViT-S/16 did not improve on SwinV2-T, RAD-DINO ViT-B/14 missed the rule on
PSNR by 0.05~dB but was continued at the full budget as an exploratory encoder, and ConvNeXt-V2-T met it.

\emph{Ours} is the ensemble of three ConvNeXt-V2-T UNet++/CBAM suppressors (seeds 0--2) trained on the Qure.ai projections with twin weight
$w=16$. It replaced the planned SwinV2-T configuration under a replacement rule declared before the final runs; the screening rule only
decided which encoders were trained at the full budget, and the replacement rule applies to every full-budget candidate. At two seeds
(development data) RAD-DINO reached $A=0.046$ and bone-region PSNR$_{\mathrm{cal}}=21.92$\,dB and ConvNeXt-V2-T $A=0.050$ and $22.01$\,dB;
both met the replacement rule, and the tie-break declared in advance (mean absorption within 0.005, then the higher bone-region PSNR) chose
ConvNeXt-V2-T, the only encoder that passed both rules, so excluding RAD-DINO would not have changed the choice. With the third seeds the
absorption gap is 0.0055, just outside the band; the choice was not reopened, and both encoders are reported. The SwinV2-T and RAD-DINO
pairs (twin and no twin, three seeds each) are secondary.

\subsubsection{Comparators}
Comparators are the \emph{no-twin control} (same encoder, data and recipe, $w=0$, three seeds); the v1 model trained on LIDC-IDRI only
(\texttt{clia}, the planned pick of the v1 study); the publicly released weights of Angaitkar et al. \cite{angaitkar2026anatomy} (SwinV2-B
at $1024$~px; research licence), called the \emph{reference model}; and three controls that need no training: identity, CLAHE
\cite{zuiderveld1994clahe} and a homomorphic filter \cite{oppenheim1968homomorphic}. All methods are scored with the same code on the same
images. BS-LDM \cite{sun2026bsldm}, trained on dual-energy pairs of the SZCH hospital, is shown for orientation only.

\subsubsection{Deep ensemble, split-conformal intervals and recalibration}
\label{sec:m-ensemble}
The reported model is a deep ensemble of $K=3$ suppressors trained with seeds $0$, $1$ and $2$ \cite{lakshminarayanan2017ensembles}; seed
$3$ is a reserve. The bone image is the ensemble mean, the soft image is $x$ minus that mean, and the per-pixel spread is the ensemble
standard deviation:
\begin{equation}
\begin{gathered}
\bar f_p=\frac1K\sum_{k=1}^{K}f_k(x)_p,\qquad \hat s=x-\bar f,\\
\sigma_p=\Big(\frac{1}{K-1}\sum_{k=1}^{K}\big(f_k(x)_p-\bar f_p\big)^2\Big)^{1/2}.
\end{gathered}
\label{eq:ensemble}
\end{equation}
With the score $r_p=|b_p-\bar f_p|/(\sigma_p+\varepsilon)$, $\varepsilon=10^{-3}$, a split-conformal scale $q$ is the
$\lceil (n+1)(1-\alpha)\rceil$-th smallest of the $n$ scores pooled over calibration pixels of the Qure.ai validation CTs, and the
interval
\begin{equation}
\mathcal{C}_p=\big[\bar f_p-q\,(\sigma_p+\varepsilon),\;\bar f_p+q\,(\sigma_p+\varepsilon)\big]
\label{eq:conformal}
\end{equation}
has nominal coverage $1-\alpha=0.9$ \cite{lei2018conformal,angelopoulos2023conformal}. Calibration treats pixels as exchangeable, so the
guarantee is marginal over pixels, not over images.

\emph{Recalibration on real pairs (H9, post hoc).} After the intervals were found not to transfer to real radiographs
(Section~\ref{sec:res-uq}), an exploratory recalibration on the 24 SZCH development pairs was added and tested on the 51 held-out pairs
(pseudo-truth bone $b=x-s_{\mathrm{DES}}$ after a gain and offset fit). Each pixel is assigned a stratum
\begin{equation}
\begin{gathered}
k(p)=3\cdot\mathbf{1}\big[\bar f_p>0.25\,p_{99}(\bar f)\big]+t(\sigma_p)\in\{0,\dots,5\},\\
q_k=r^{(k)}_{(\lceil (n_k+1)(1-\alpha)\rceil)},
\end{gathered}
\label{eq:mondrian}
\end{equation}
where $t(\sigma_p)\in\{0,1,2\}$ is the spread tercile (edges from the calibration pixels) and $r^{(k)}_{(j)}$ the $j$-th smallest of the
$n_k$ calibration scores in stratum $k$; the interval of Eq.~\eqref{eq:conformal} then uses $q_{k(p)}$ (Mondrian calibration). The global
variant uses one stratum, and a per-image variant replaces the pooled pixel scores by the per-image $(1-\alpha)$ quantiles, calibrated at
level 0.1. The sparsification curve is the mean absolute error of the pixels that remain after removing the fraction $f\in[0,0.95]$ with
the largest spread; AUSE is the area between this curve and the oracle curve that removes pixels by their error (0 is perfect).

\subsubsection{Audit sites and lesions}
\label{sec:m-audit}
Lesions come from the CTs of the LIDC-IDRI test split only, so no audited lesion was seen in training. On each real radiograph six sites are
drawn at random positions inside the lung fields at $\alpha=1$; the sites depend only on the image, its lung mask, the lesion library and a
fixed seed, and are shared by every model. Lung masks come from the PSPNet anatomical segmentation model of TorchXRayVision
\cite{cohen2022torchxrayvision}, eroded by 49~px at $1024$~px so that a site lies inside the lung field; the model comes from the ChestX-Det
project \cite{lian2021chestxdet}, whose paper does not fully document its training data. The library holds 296 nodules (4.0 to 37.6~mm,
median 7.1~mm); six sites per image give 1,470 sites on VinDr-RibCXR, 1,482 on JSRT and 450 on SZCH. Because the PSPNet masks differ
slightly between GPUs, the masks were exported once on one machine and every evaluation reads them, so every model is scored on identical
sites.

The bone overlap of a site is measured with independent information, never with a model output: the radiologist rib masks of VinDr-RibCXR
\cite{nguyen2021ribcxr} for the confirmatory audit, and the real dual-energy bone signal on SZCH. A site \emph{overlaps bone} when at least
15\,\% of its support lies on an annotated rib. On JSRT no bone reference exists; the protocol named the consensus of the bone maps of ours
and of the reference model, but the audit took the overlap from the reference model's bone map alone (a secondary analysis;
Supplementary Section~S1).

\emph{Why the audit is not circular.} The twin regulariser pastes lesions to teach invariance, and the audit pastes lesions to measure it,
so a model could pass the audit by learning the insertion artefact. Four design choices limit this. The audit lesions come from CTs that
were never used for training; the audit images are real radiographs, whereas the training composites are synthetic; the confirmatory audit
uses rib masks that no model has seen; and the claim that the regulariser helps is also tested with evidence that does not use insertion:
contrast of real nodules on JSRT and Node21, detection of real lesions, and detection recall grouped by bone overlap (H4$'$). These choices
limit the circularity but cannot remove it: a model that recognises pasted signal in general could still score well, which is why the
evidence without insertion is reported beside the audit (Section~\ref{sec:limitations}).

\emph{Audit variants.} On the Qure.ai test split the lesion is added to the soft component, the true bone image is known, and the bone
error inside the lesion support can be measured against it in addition to $A$ (exact-truth audit, Supplementary Table~S5). The SZCH audit
ran on all 75 SZCH images (450 sites) during development and is reported as development evidence only; the SZCH contrasts in the result
tables use its 51 held-out images (306 sites). In an amplitude sweep, $A$ is recomputed for $\alpha\in\{0.25,0.5,1,2,4\}$ on the 450 SZCH
sites (development data, Supplementary Table~S6). For $\alpha\le1$ the absorption is flat (v1 model 0.301--0.306, no-twin model
0.215--0.218, reference model 0.134--0.136), so the leak is linear in the lesion amplitude; at $\alpha=2$ and 4 it falls for every model,
where the added signal starts to saturate at the top of the intensity range.

\subsubsection{Fidelity and suppression measures on real images}
\label{sec:m-fidelity}
\emph{Paired dual-energy data.} Tone differs between a radiograph and its DES image, so metrics on the SZCH pairs are reported as is
(``raw'') and after a per-image least-squares gain and offset of the estimate $\hat s$ onto the reference $s$ (``calibrated'', primary):
\begin{equation}
\begin{gathered}
(\hat a,\hat c)=\arg\min_{a,c}\sum_p\big(a\,\hat s_p+c-s_p\big)^2,\\
\mathrm{PSNR}_{\mathrm{cal}}=10\log_{10}\frac{1}{\frac{1}{|\Omega|}\sum_{p\in\Omega}\big(\mathrm{clip}_{[0,1]}(\hat a\,\hat s_p+\hat c)-s_p\big)^2},
\end{gathered}
\label{eq:psnrcal}
\end{equation}
with intensities in $[0,1]$ and $\Omega$ the evaluated region; SSIM$_{\mathrm{cal}}$ uses the same calibrated estimate. Calibrated PSNR and
SSIM are also reported inside the region where DES removed the most signal (top 10\,\% of calibrated radiograph minus DES, the bone region)
and outside it. Paired fidelity also uses PSNR and SSIM \cite{wang2004ssim} on the whole image, band-limited SSIM, MS-SSIM
\cite{wang2003msssim}, LPIPS \cite{zhang2018lpips} and a high-frequency energy ratio; the same smoothing is applied to every method and never
tuned.

\emph{Rib benchmark.} On VinDr-RibCXR (245 images; evaluated locally, never used for training) we compare the gradient magnitude
$G(I)=\lvert\nabla(g_{1}*I)\rvert$ (Sobel operator after a Gaussian blur of 1~px) of the suppressed image $\hat s$ with that of the
radiograph $x$:
\begin{equation}
\begin{gathered}
R_{\mathrm{rib}}=1-\frac{\sum_{p\in\mathcal{R}}G(\hat s)_p}{\sum_{p\in\mathcal{R}}G(x)_p},\\
R_{\mathrm{non}}=1-\frac{\sum_{p\in\mathcal{N}}G(\hat s)_p}{\sum_{p\in\mathcal{N}}G(x)_p},
\end{gathered}
\label{eq:rib}
\end{equation}
where $\mathcal{R}$ is the rib-boundary band inside the lungs and $\mathcal{N}$ the lung pixels away from ribs. We report both, their
difference $R_{\mathrm{rib}}-R_{\mathrm{non}}$ (selectivity) and the ROC-AUC of the gradient-reduction map for separating the two sets of
pixels.

\emph{Real-nodule contrast.} For a nodule box $\mathcal{B}$ on JSRT or Node21 and the ring $\mathcal{U}$ around it (the box grown by half
its size on each side),
\begin{equation}
\begin{gathered}
\mathrm{CNR}(I)=\frac{\mathrm{mean}_{\mathcal{B}}(I)-\mathrm{mean}_{\mathcal{U}}(I)}{\mathrm{sd}_{\mathcal{U}}(I)},\\
\Delta\mathrm{CNR}=\mathrm{CNR}(\hat s)-\mathrm{CNR}(x),
\end{gathered}
\label{eq:cnr}
\end{equation}
compared between suppressors as paired differences per nodule.

\subsubsection{Site-appearance calibration (exploratory)}
\label{sec:m-sitecal}
BS-LDM, trained on dual-energy pairs of the SZCH hospital, scores far higher than any CT-supervised model on the held-out pairs
(Section~\ref{sec:res-fidelity}). After the H2$'$ results were known, diagnostics on the 24 development pairs only asked why. Subtracting
more of the bone estimate lowered PSNR$_{\mathrm{cal}}$ rather than raising it; a bone signal known exactly but predicted at 512\,px would
reach 41.7\,dB, so resolution does not limit it; and most of the error was broader than 64\,px and present in the unprocessed radiograph
as well. The dual-energy images of this site therefore differ from a radiograph minus its bone image by a smooth appearance field that the
per-image gain and offset of the metric cannot absorb (Fig.~\ref{fig:sitecal_diag}). A protocol written before any held-out pair was scored
with it tests a smooth correction of this field:
\begin{equation}
\begin{gathered}
\hat s'=\hat s+U(F w),\\
w=\arg\min_{w}\sum_{i\in\mathcal{D}}\big\lVert \Pi\,G_{16}(s_i-\hat s_i)-F_i w\big\rVert_2^2+\lambda N\lVert w\rVert_2^2,
\end{gathered}
\label{eq:sitecal}
\end{equation}
where $\mathcal{D}$ are the 24 development pairs with dual-energy soft-tissue images $s_i$ and estimates $\hat s_i$, $G_{16}$ is a Gaussian
blur of 16\,px at 1024\,px, $\Pi$ samples a $64\times64$ grid, $F_i$ holds on that grid the features $[1,u,v,u^2,v^2,uv,I,I^2,B,Iv,Iu]$
(position $u,v$; blurred radiograph $I$ and bone estimate $B=x-\hat s$), $\lambda=10^{-3}$, $N$ is the number of grid samples and $U$ is
bicubic upsampling. One $w$ is fitted per method in the same way (identity, ours, no-twin control, v1 model, reference model) and applied
unchanged to the 51 held-out pairs; BS-LDM, trained on this site's images, is not corrected. The fitted weight on $B$ is about $+0.35$ for
every CT-supervised model, so the correction adds back about a third of the smoothed bone estimate. The correction adds only a smooth field
and so cannot restore an erased lesion. Its hypotheses H16--H18 are stated in Section~\ref{sec:m-stats}. The analysis is exploratory: it
re-uses held-out pairs already used for H2$'$, H2$'$ stays as reported, and calibrated numbers are always shown beside the uncorrected ones.

\begin{figure*}[!t]
\centering
\safefig[0.9\textwidth]{figures/fig_sitecal_diag.pdf}
\caption{Development-pair diagnostics behind the site calibration (24 SZCH development pairs; the held-out pairs were not used). (a)
Calibrated PSNR of the whole image and of the bone region when the bone estimate is subtracted with gain $k$ ($k=1$ as trained):
subtracting more lowers PSNR for both CT-supervised models. (b) Whole-image PSNR of ours as is, after the leave-one-out site correction of
Eq.~\eqref{eq:sitecal}, if every error component broader than 64 or 16\,px were removed, and for an oracle that knows the dual-energy bone
signal exactly but only at 512\,px.}
\label{fig:sitecal_diag}
\end{figure*}

% ---------------------------------------------------------------------------------------------------------------------
\subsection{Downstream Detection and Classification}
\label{sec:m-downstream}
The downstream classifier is a DenseNet-121 \cite{huang2017densenet} and the detector a Faster R-CNN with a ResNet-50 FPN backbone
\cite{ren2017fasterrcnn}, with identical settings for every input variant. The classifier sees $256$~px images, is trained for 10 epochs
with AdamW \cite{loshchilov2019adamw} (learning rate $3\times10^{-4}$, weight decay $10^{-4}$, batch 64, cosine schedule) and keeps the
epoch with the best validation macro-AUROC. The detector sees $512$~px images, is trained for 12 epochs with SGD (learning rate 0.01,
momentum 0.9, weight decay $10^{-4}$, batch 8) and keeps the epoch with the best validation CPM. Each input is trained with three seeds:
the original radiograph, the soft-tissue image of ours or of the reference model, CLAHE, and a multi-channel input that stacks images as
channels (for the classifier the original and our soft image, for the detector the original, our soft and our bone image).
Classification uses NIH ChestX-ray14 (macro-AUROC). Detection uses Node21, TBX11K and VinBigData and is scored by FROC sensitivity at fixed false-positive rates, averaged as
the competition performance metric (CPM) \cite{vanginneken2010anode09,niemeijer2011cpm}, and by mAP@50; detection metrics treat tied scores
as one threshold, so they do not depend on the order in which a sort returns ties.

% ---------------------------------------------------------------------------------------------------------------------
\subsection{Study~II: Multi-Stream Fusion with Capacity-Matched Controls}
\label{sec:m-study2}

\subsubsection{Rationale}
With the frozen suppressor ensemble $\bar f$ of Study~I, a classifier can receive the radiograph $X$, the soft-tissue image
$S=X-\bar f(X)$ and the bone image $\hat B=\mathrm{clip}_{[0,1]}(X-S)$ on the stored 8-bit scale. Because $(S,\hat B)$ is a deterministic
function of $X$, the inputs carry no new information about a label $Y$:
\begin{equation}
I\big(Y;\,X,S,\hat B\big)=I\big(Y;\,X\big).
\label{eq:info}
\end{equation}
A gain can therefore come only from the prior the suppressor learned from CT, or from the inductive bias of the network that reads the
extra images; a network that reads three images also spends three times the encoder computation. The study separates these sources with a
control of the same capacity, and it analyses the effect per finding group, because suppression could help focal lesions behind ribs while
removing cues for diffuse or pleural findings.

\subsubsection{Input configurations and network}
Five configurations were trained with identical settings (Figs.~\ref{fig:architecture}b and~\ref{fig:cls_design}):
\emph{orig} ($X$), \emph{soft} ($S$), \emph{early} ($X$, $S$ and $\hat B$ as the three input channels of one stream), \emph{hybrid} (three
streams $X$, $S$, $\hat B$) and \emph{ctrl} (the hybrid network with $X$ in all three streams). The soft-tissue image comes from our
suppressor ensemble run at its input size of 512\,px; classifiers see all images at $256\times256$\,px. One DenseNet-121 encoder $E$
\cite{huang2017densenet} with shared weights, pretrained on ImageNet \cite{deng2009imagenet,wightman2019timm}, processes each stream $u_s$.
Each stream is standardised per image and passed through its own learned affine map, because the bone image spans a narrow intensity range;
the $8\times8$ encoder output is projected to $D=256$-dimensional tokens with a stream embedding $e_s$ and a two-dimensional sine position
code $\pi_n$:
\begin{equation}
\begin{gathered}
\tilde u_s=\gamma_s\,\frac{u_s-\mathrm{mean}(u_s)}{\mathrm{sd}(u_s)}+\beta_s,\\
z_{s,n}=W\,E(\tilde u_s)_n+e_s+\pi_n\in\mathbb{R}^{D}.
\end{gathered}
\label{eq:tokens}
\end{equation}
A two-layer decoder with eight heads and one learned query $q_c$ per class attends to the tokens of all streams, as in query-based
multi-label heads \cite{liu2021query2label,ridnik2023mldecoder}, and a linear map per class gives the logit $\ell_c$ from the decoder
output $o_c$. With keys $k_{s,n}$, the head-averaged cross-attention of the last layer defines the routing of class $c$ to stream $s$:
\begin{equation}
\begin{gathered}
a_{c,(s,n)}=\operatorname*{softmax}_{(s,n)}\Big(\frac{q_c^{\top}k_{s,n}}{\sqrt{D}}\Big),\\
\rho_{c,s}=\sum_{n}a_{c,(s,n)},\qquad
\ell_c=w_c^{\top}o_c+b_c.
\end{gathered}
\label{eq:routing}
\end{equation}
Single-stream configurations use the same encoder and decoder over 64 tokens. The shared encoder gives the hybrid the parameters of a
single-stream model plus the decoder, at three times the encoder computation; \emph{ctrl} has exactly the hybrid's parameters and
computation but no suppression.

\begin{figure*}[!t]
\centering
\safefig[0.9\textwidth]{figures/fig_study_design.pdf}
\caption{Fusion study design. (a) The three inputs for an NIH ChestX-ray14 test image: the radiograph $X$, the soft-tissue image $S$
predicted by the suppressor and the bone image $\hat B=X-S$ (box: expert annotation of a focal finding). (b) The five configurations, all
with the same DenseNet-121 encoder, class-query decoder and training settings. In the hybrid and the control one encoder with shared
weights processes each stream and the decoder attends to the tokens of all three; the control receives $X$ three times. Bottom: the
contrasts tested.}
\label{fig:cls_design}
\end{figure*}

\subsubsection{Training}
The loss is the binary cross-entropy over the annotated labels $\mathcal{V}$ of a batch (one-hot targets for the three BDCXR-3257 classes):
\begin{equation}
\mathcal{L}_{\mathrm{cls}}=-\frac{1}{|\mathcal{V}|}\sum_{(i,c)\in\mathcal{V}}\Big[y_{ic}\log\sigma(\ell_{ic})+(1-y_{ic})\log\big(1-\sigma(\ell_{ic})\big)\Big].
\label{eq:bce}
\end{equation}
Training used AdamW \cite{loshchilov2019adamw} (learning rate $3\times10^{-4}$, weight decay $10^{-4}$), batch size 32, 20 epochs with a
cosine schedule after 200 warm-up steps, and mixed precision. Augmentation: random rotation by up to $7^\circ$, scaling by 0.9--1.1,
translation by up to 3\,\% of the image size, and brightness and contrast changes, identical across the channels of an image. The epoch with
the best validation macro-AUROC was kept. Every configuration was trained with seeds 0--4 on NIH ChestX-ray14 and VinBigData and with seeds
0--4 in each of the five BDCXR-3257 folds of Section~\ref{sec:m-bdcxr}, giving 175 runs. A configuration's prediction is the seed ensemble
\begin{equation}
\bar p_c(x)=\frac15\sum_{k=0}^{4}\sigma\big(\ell^{(k)}_c(x)\big),
\label{eq:seedens}
\end{equation}
and all tests compare seed ensembles. Runs used cloud T4 GPUs; a run longer than one session resumed from an end-of-epoch checkpoint, which
gave the same predictions as uninterrupted training in a test.

\subsubsection{Encoder selection}
We screened four encoders on the NIH validation split (configuration \emph{orig}, seed 0, 10 epochs): DenseNet-121, ConvNeXt-V2-T
\cite{woo2023convnextv2}, RAD-DINO \cite{perezgarcia2025raddino} and SwinV2-B \cite{liu2022swinv2}. The best encoder would go forward with
DenseNet-121, and an encoder within 0.005 of the best would yield to a cheaper one. All four reached validation macro-AUROC between 0.8384
and 0.8430, and DenseNet-121 was the cheapest, so it went forward alone. ConvNeXt-V2-T failed to train at the planned learning rate;
following the rule of the suppressor study, the three encoders without batch normalisation were screened again with the encoder learning
rate reduced fourfold, a change decided before the other screen results were seen (Supplementary Section~S1).

\subsubsection{Finding groups}
Before any run, the classes were grouped as \emph{focal} (NIH: nodule, mass; VinBigData: nodule/mass), \emph{diffuse} (effusion,
infiltration, consolidation, oedema, pneumonia; VinBigData: pleural effusion, infiltration, consolidation, lung opacity, interstitial lung
disease) and \emph{boundary} (cardiomegaly, atelectasis, pneumothorax, pleural thickening, emphysema, fibrosis, hernia; VinBigData:
cardiomegaly, atelectasis, pneumothorax, pleural thickening, aortic enlargement, pulmonary fibrosis, calcification). VinBigData's ``other
lesion'' belongs to no group, and BDCXR-3257 (normal, bacterial or viral pneumonia) enters only the overall tests. For a group $G$ of classes
and two configurations $m$ and $m'$,
\begin{equation}
\begin{gathered}
\mathrm{AUROC}_{\mathrm{macro}}(m)=\frac1C\sum_{c=1}^{C}\mathrm{AUROC}_c(m),\\
\Delta_G(m,m')=\frac{1}{|G|}\sum_{c\in G}\big(\mathrm{AUROC}_c(m)-\mathrm{AUROC}_c(m')\big).
\end{gathered}
\label{eq:group}
\end{equation}

\subsubsection{Attribution maps and the pointing game}
\label{sec:m-xai}
For each hybrid model, Grad-CAM++ \cite{chattopadhay2018gradcampp} was computed per stream on the last encoder block. Integrated Gradients
\cite{sundararajan2017axiomatic} was computed on the three inputs jointly against a blurred baseline ($\sigma=8$\,px, 32 steps), because a
black image has no contrast after standardisation; the decoder attention was read from its last layer. A target $t$ is a pair of test image
and box class; a map $M_m$ hits when its maximum lies inside a box of that class, with a tolerance of 15\,px at 256\,px fixed in advance
\cite{zhang2018topdown}. With $h^{(k)}_m(t)\in\{0,1\}$ the hit of map $m$ of the seed-$k$ hybrid model, focal targets $\mathcal{F}$
(masses and nodules) and the other targets $\mathcal{O}$ (six box classes),
\begin{equation}
\begin{gathered}
d(t)=\frac15\sum_{k=0}^{4}\big(h^{(k)}_{S}(t)-h^{(k)}_{X}(t)\big),\\
I=\frac{1}{|\mathcal{F}|}\sum_{t\in\mathcal{F}}d(t)-\frac{1}{|\mathcal{O}|}\sum_{t\in\mathcal{O}}d(t),
\end{gathered}
\label{eq:pointing}
\end{equation}
and H10 tests $I>0$ with an image-cluster bootstrap (\num{10000} resamples; Eq.~\eqref{eq:bootp}).

\subsubsection{Classification uncertainty}
\label{sec:m-clsuq}
\emph{NIH: false-negative risk control.} An image's label set holds every class with ensemble probability at least $1-\lambda$; the threshold
is calibrated on the validation split (\num{8634} images) by conformal risk control \cite{angelopoulos2022crc} and the false-negative rate is
measured on the test split:
\begin{equation}
\begin{gathered}
\mathcal{C}_\lambda(x)=\{c:\bar p_c(x)\ge1-\lambda\},\\
\hat R_n(\lambda)=\frac1n\sum_{i=1}^{n}\Big(1-\frac{|\mathcal{C}_\lambda(x_i)\cap Y_i|}{|Y_i|}\Big),\\
\hat\lambda=\min\Big\{\lambda:\frac{n\,\hat R_n(\lambda)+1}{n+1}\le\alpha\Big\},
\end{gathered}
\label{eq:crc}
\end{equation}
over the $n$ calibration images with at least one positive label $Y_i$, at $\alpha=0.10$.

\emph{BDCXR-3257: class-conditional sets.} With probabilities normalised to sum to one, $\hat p_y=\bar p_y/\sum_{y'}\bar p_{y'}$, the score
$1-\hat p_y$ is calibrated separately for each true class \cite{vovk2005algorithmic}, on the other folds' out-of-fold predictions for each
test fold:
\begin{equation}
\begin{gathered}
\mathcal{C}(x)=\{y:1-\hat p_y(x)\le q_y\},\\
q_y=\text{the }j_y\text{-th smallest of }\{1-\hat p_{y_i}(x_i):y_i=y\},\\
j_y=\lceil (n_y+1)(1-\alpha)\rceil.
\end{gathered}
\label{eq:mondriansets}
\end{equation}

\emph{Selective classification.} Images are ranked by a confidence score and the most confident fraction is kept
\cite{geifman2017selective,geifman2019bias}. Score A is one minus the entropy $H$ of the seed ensemble (mean binary entropy over the 14
labels for NIH, Shannon entropy of the normalised probabilities for BDCXR-3257); score B adds the suppressor's own uncertainty,
$-(z_H+z_\sigma)$, where $\sigma$ is the spread of Eq.~\eqref{eq:ensemble} averaged over the lung field and $z$ standardises over the
evaluated images of a data set. With the losses $\ell_{(1)},\dots,\ell_{(N)}$ sorted by decreasing confidence (Brier score for NIH, top-1
error for BDCXR-3257),
\begin{equation}
\mathrm{AURC}=\frac1N\sum_{k=1}^{N}\frac1k\sum_{j=1}^{k}\ell_{(j)}.
\label{eq:aurc}
\end{equation}

% ---------------------------------------------------------------------------------------------------------------------
\subsection{External Validation on Paediatric Radiographs}
\label{sec:m-bdcxr}
The protocol of this validation was fixed before any model of this study saw a BDCXR-3257 image; its one later change, the duplicate
grouping, was made from the images alone before any model ran (Supplementary Section~S1). All suppressors are the frozen models; only the
downstream classifiers are trained. (E1, primary) Pneumonia classification with three one-vs-rest outputs, inputs as in H4$'$b (original,
ours, reference model, CLAHE, original plus ours), the classifier settings of the downstream protocol, five folds stratified by class over
the near-duplicate groups (fold seed 0) with 10\,\% of each training part, by whole groups, for epoch selection, and three seeds per fold and
input (75 runs). Out-of-fold probabilities, averaged over seeds, give one prediction per image. The primary measure is the pooled
macro-AUROC; ours against the original image is non-inferior if the lower bound of the 95\,\% interval of the difference (group bootstrap)
exceeds $-0.01$, and superior if it exceeds 0; per class, paired DeLong tests with Holm over the three classes. (E2, secondary) The lesion
audit on a fixed class-stratified subset of 300 images (seed 7), six sites per image, with bone overlap from the reference model's bone map;
images whose lung mask covers less than 2\,\% of the image are skipped and counted (none was). (E3, descriptive) The per-image ensemble
spread over the lung field, compared with the held-out SZCH pairs and correlated with age.

% ---------------------------------------------------------------------------------------------------------------------
\subsection{Statistical Analysis}
\label{sec:m-stats}

\subsubsection{General framework}
Hypotheses, margins, test sets and tests were written down before the data they govern were evaluated, and they carry the labels under
which they were specified (H1$'$--H5$'$ for the final suppressors, H1 for the v1 models, H6--H10 for the fusion study, H16--H18 for the site
calibration and E1--E3 for the external validation). Analyses added later are labelled exploratory or post hoc, every change made after the
analyses started is listed in Supplementary Section~S1, and all primary and secondary results are reported whatever their direction.

All comparisons are paired. Confidence intervals are percentile bootstrap intervals \cite{efron1993bootstrap} (\num{10000} resamples unless
stated) over patients when several units come from one patient; paired differences are tested with the Wilcoxon signed-rank test
\cite{wilcoxon1945individual}; AUCs are compared with the paired DeLong test \cite{delong1988auc}. Seeds are averaged within unit before
testing and their spread is reported. Multiplicity is handled by the Holm procedure \cite{holm1979simple} within each family: with the $m$
$p$-values of a family ordered as $p_{(1)}\le\dots\le p_{(m)}$,
\begin{equation}
\tilde p_{(j)}=\max_{i\le j}\,\min\big\{1,\,(m-i+1)\,p_{(i)}\big\}.
\label{eq:holm}
\end{equation}
Confidence intervals are unadjusted 95\,\% intervals; effect sizes are always reported with them, and a hypothesis that fails is reported as
failed. Non-inferiority uses the paired mean difference $\bar\Delta=\frac1n\sum_i(m_i^{\mathrm{ours}}-m_i^{\mathrm{ref}})$ and the
one-sided 95\,\% lower bound $L_{95}$ of its image-level bootstrap distribution; it holds if
\begin{equation}
L_{95}(\bar\Delta)>-\delta
\label{eq:ni}
\end{equation}
for a margin $\delta$ fixed in advance.

\subsubsection{Study~I hypotheses (H1$'$--H5$'$)}
\begin{description}
\item[H1$'$, dose-response.] On VinDr-RibCXR (245 images, six sites each, lesions from LIDC test CTs, $\alpha=1$), the per-site absorption
 $A$ averaged over seeds falls with the twin weight. With $\bar A_{s,w}$ the absorption at site $s$ averaged over the seeds trained with
 twin weight $w$, the slope $\beta_1$ of
 \begin{equation}
 \bar A_{s,w}=\beta_0+\beta_1\log_2(1+w)+e_{s,w}
 \label{eq:dose}
 \end{equation}
 has an image-cluster bootstrap interval, and H1$'$ passes if its upper bound is below zero. The same analysis on the Qure.ai test split
 (exact truth) is secondary.
\item[H2$'$, fidelity.] On the 51 held-out SZCH pairs, bone-region PSNR of ours against the reference model is non-inferior
 (Eq.~\eqref{eq:ni}) with margin $\delta=0.5$~dB; SSIM has a margin of 0.01 and is reported either way; ours against the no-twin control
 uses a margin of 0.25~dB (the cost of the regulariser). Metrics are computed calibrated (primary) and raw (Eq.~\eqref{eq:psnrcal}).
\item[H3$'$, lesion absorption.] On VinDr-RibCXR sites, $A$(ours)$<A$(reference model) and $A$(ours)$<A$(no-twin control), as paired
 site-level differences with image-cluster bootstrap intervals, Holm over the two contrasts, and an interaction between sites that overlap
 an annotated rib and sites that do not.
\item[H4$'$, evidence without insertion.] (a) $\Delta$CNR of real nodules on JSRT and Node21 (Eq.~\eqref{eq:cnr}), ours against the
 reference model and against the no-twin control, for all lesions and for lesions with at least 15\,\% bone overlap. (b) Downstream: NIH
 ChestX-ray14 macro-AUROC per input (original, ours, reference model, CLAHE, original plus ours; three seeds); TBX11K, Node21 and
 VinBigData detection by FROC and mAP@50; recall at two false positives per image for lesions with and without bone overlap, as a
 difference in differences against the original input.
\item[H5$'$, reliability.] Spearman correlation between ensemble spread and absolute bone error on the SZCH held-out pairs (pseudo-truth
 bone from the least-squares-aligned DES image); split-conformal coverage calibrated on Qure.ai validation CTs and measured on the Qure.ai
 test split (nominal 90\,\%) and, as measured under shift, on SZCH.
\end{description}
Site-level differences between models are paired, with image-cluster bootstrap intervals. Bootstrap intervals use \num{10000} resamples for
mean absorption and contrasts, \num{2000} for the overlap interaction and 500 for the dose-response slope, always resampling whole images
with replacement. The hypothesis H1 of the v1 models (insertion lowers absorption on synthetic LIDC test CTs at non-inferior bone error) is
tested in the same way on exact truth.

\emph{Site calibration (exploratory, Section~\ref{sec:m-sitecal}).} H16 (primary) is a gain in the bone-region PSNR$_{\mathrm{cal}}$ of ours
from the calibration (one-sided 95\,\% lower bound of the paired difference above 0; image bootstrap, \num{10000} resamples) and H16b the
same for whole-image PSNR$_{\mathrm{cal}}$. H17 is the non-inferiority of the calibrated ours to the calibrated reference model in
bone-region PSNR$_{\mathrm{cal}}$ (margin 0.5\,dB, as H2$'$). H18 is a lower absorption of the calibrated ours than of the calibrated
reference model on VinDr-RibCXR, from one audit run with all four models so that they share every site (image-cluster bootstrap, two-sided
95\,\% upper bound below 0). The rib benchmark of the calibrated models is descriptive. H16 is the only confirmatory test of this analysis;
the others are reported without multiplicity correction.

\subsubsection{Study~II hypotheses (H6--H10)}
H9 concerns the suppressor's intervals (Section~\ref{sec:m-ensemble}); the others concern the classifiers.
\begin{description}
\item[H6 (primary).] The hybrid has a higher macro-AUROC than \emph{orig}; one test per data set, Holm correction over the three.
\item[H7 (secondary).] Hybrid against \emph{ctrl} (does a gain come from the suppression or from the extra computation?) and against
 \emph{early} (do separate streams beat channel stacking?); one Holm family over the six tests.
\item[H8 (secondary).] $\Delta_G$ of Eq.~\eqref{eq:group} for hybrid and for \emph{soft} against \emph{orig} in each finding group; Holm
 within data set over the six tests. Per class, hybrid against \emph{orig} and against \emph{early} with the paired DeLong test
 \cite{delong1988auc}, Holm within data set over classes and both contrasts.
\item[H10 (secondary).] In the hybrid models, the pointing-game hit rate of the soft-tissue stream minus that of the radiograph stream is
 larger for focal targets than for the other targets (Eq.~\eqref{eq:pointing}).
\end{description}
Every contrast $\Delta$ of two seed ensembles is resampled with a paired cluster bootstrap: replicate $b$ draws clusters with replacement
(patients for NIH, images for VinBigData, near-duplicate groups for BDCXR-3257), recomputes $\Delta^{*}_b$ with the same resamples for every
contrast of a data set, and with $B=\num{10000}$ replicates gives the percentile interval $[\Delta^{*}_{(0.025)},\Delta^{*}_{(0.975)}]$ and
\begin{equation}
\begin{gathered}
p=\min\big\{1,\;2\min(p_{-},p_{+})\big\},\\
p_{-}=\frac{1+\#\{b:\Delta^{*}_b\le0\}}{B+1},\quad p_{+}=\frac{1+\#\{b:\Delta^{*}_b\ge0\}}{B+1}.
\end{gathered}
\label{eq:bootp}
\end{equation}
AUPRC differences are reported with bootstrap intervals only. The hypotheses, the code and the split files of this study were fixed before
the encoder screen and before any classifier of this study saw a test split.
